%! TEX root = main.tex

\chapter{Homework 1}

\section{Weak Form and FE Discretization}
\label{sc:11}

In this chapter, the weak formulation and the corresponding finite element (FE) discretization of the boundary value problem introduced in Homework 1 are derived.

\subsection{Strong Formulation}

The one-dimensional wave propagation problem is considered:

\begin{equation}
\left\{
\begin{aligned}
&\rho(x) \omega^{2} u(x) + \frac{\partial}{\partial x} \left(\mu(x) \frac{\partial u(x)}{\partial x}\right) = f \quad x \in (0,L)\\
&\mu(0) \left.\frac{\partial u(x)}{\partial x}\right|_{x = 0} = g_N \\
&u(L) - \sqrt{\frac{\mu(L)}{\rho(L)}} \frac{i}{\omega} \left.\frac{\partial u(x)}{\partial x}\right|_{x = L} = g_A
\end{aligned}
\right.
\label{eq:strongform}
\end{equation}

In this report, the following notational simplifications are adopted:

\begin{itemize}
    \item \( u = u(x) \) and, analogously, for all functions depending on time;
    \item \( u' = u'(x) = \frac{\partial u(x)}{\partial x} \) and, analogously, for all spatial derivatives;
    \item \( u'(0) = \left.\frac{\partial u(x)}{\partial x}\right|_{x = 0} \) and, analogously, for derivatives calculated in \( x = L\).
\end{itemize}

The second boundary condition is an impedance (or absorbing) condition. It is useful to rewrite it in order to obtain \( u'(L) \) explicitly. From the relation above, it follows that:
\[
u'(L) = -i \omega \sqrt{\frac{\rho(L)}{\mu(L)}} \left( u(L) - g_A \right)
\]
For convenience, the impedance coefficient is introduced as
\[
\beta = \sqrt{\rho(L)\mu(L)}
\]

\subsection{Weak Formulation}

To obtain the weak formulation, the strong form of the equation is multiplied by a test function \( v(x) \in H^1(0,L) \) and integrated over \( (0,L) \):
\[
\int_0^L \rho \omega^2 u v dx + \int_0^L \frac{d}{dx} \left( \mu u' \right) vdx = \int_0^L f v dx
\]
The second term is then integrated by parts:
\[
\int_0^L \frac{d}{dx}(\mu u') v dx = \left[\mu u' v\right]_{0}^{L} - \int_0^L \mu u' v' dx
\]
Substituting into the previous expression yields:
\[
\omega^2\int_0^L \rho u v dx - \int_0^L \mu u' v' dx + \left[\mu u' v\right]_{0}^{L} = \int_0^L f vdx
\]
The boundary contributions are incorporated using the boundary conditions. At \( x=0 \), it holds that:
\[
\mu(0) u'(0) v(0) = g_N v(0)
\quad \Longrightarrow \quad 
-\mu(0) u'(0) v(0) = -g_N v(0)
\]
At \( x = L \), using \( u'(L) \) from the impedance condition, one obtains:
\[
\mu(L) u'(L) v(L) = -i \omega \beta u(L) v(L) + i \omega \beta g_A v(L)
\]
Collecting all terms, the weak formulation becomes:

\begin{equation}
\omega^2 \int_0^L \rho u v dx - \int_0^L \mu u' v' dx - i\omega\beta u(L) v(L) = \int_0^L f v dx + g_N v(0) - i\omega \beta g_A v(L)
\label{eq:weak_form}
\end{equation}

It is convenient to identify the mass term, stiffness term and the contributions arising from the boundary conditions. The bilinear forms are defined as:
\[
\mathcal{M}(u,v) = \int_0^L \rho u v dx,
\quad
\mathcal{A}(u,v) = \int_0^L \mu u' v' dx + i \omega \beta u(L) v(L)
\]
and the linear functional as:
\[
\mathcal{F}(v) = \int_0^L f v dx + g_N v(0) - i \omega \beta g_A v(L)
\]
Thus, the weak formulation can be written compactly as:

\begin{equation}
\omega^2 \mathcal{M}(u,v) - \mathcal{A}(u,v) = \mathcal{F}(v) \quad \forall v \in H^1(0,L)
\label{eq:comptweakform}
\end{equation}

\subsection{Finite Element Discretization and Algebraic System}

A partition of the interval \( (0,L) \) into \( N \) elements is introduced and the finite dimensional space \( V_h \subset H^1(0,L) \), spanned by continuous, piecewise linear basis functions \( \{\varphi_j\}_{j=1}^N \), is considered. The solution and the test function are approximated as:
\[
u_h = u_h(x) = \sum_{j=1}^N u_j(x) \varphi_j(x), \quad v_h = v_h(x) = \varphi_i(x)
\]
Substituting \( u_h \) and \( v_h \) into the equation \eqref{eq:weak_form} yields the discrete system. The entries of the mass matrix, stiffness matrix and load vector are then defined.

\begin{itemize}
    \item Mass matrix:
    \[
    M_{ij} = \int_0^L \varphi_i \varphi_jdx
    \]
    \item Stiffness matrix and impedance contribution:
    \[
    A_{ij} = \int_0^L \mu \varphi_i' \varphi_j' dx + i \omega \beta \varphi_i(L) \varphi_j(L)
    \]
    Note that \( \varphi_i(L)\neq 0 \) is only for the basis function associated with the last node of the mesh. Therefore, the impedance term modifies only the entry \( A_{NN} \).
    \item Load vector:
    \[
    F_i = \int_0^L f \varphi_i dx + g_N \varphi_i(0) - i \omega \beta g_A \varphi_i(L)
    \]
    Here \( \varphi_i(0) \) is nonzero only for the basis function associated with the first node, while \( \varphi_i(L) \) is nonzero only for the last one.
\end{itemize}

Collecting all contributions, the linear system is obtained:

\begin{equation}
\omega^2 \mathbf{M} \underline{u} - \mathbf{A} \underline{u} = \underline{F}
\label{eq:algsystm}
\end{equation}

which corresponds to the discretized weak formulation and matches the form requested in the homework assignment. This completes the derivation of the weak formulation and its finite element discretization using linear basis functions.